\section{Certificates after bounded monomial reformulation}
\label{sec:monomial-reformulation}
The preceding transfer survives coordinate branching on products of
boundedly many variables. The number of auxiliary coordinates is
arbitrary; the count depends on the original variables covered by
independent parity restrictions.

Consider a map $h:[-1,1]^n\to[-1,1]^N$ of the form
\begin{equation}\label{eq:signed-monomial-map}
 h_j(x)=\sigma_j\chi_{A_j}(x),\qquad
 \sigma_j\in\{-1,1\},\qquad 1\le|A_j|\le D,
 \qquad j=1,\ldots,N,
\end{equation}
where $D\ge1$ is an integer. Include every original coordinate as an
identity coordinate. The nonempty supports are squarefree; repeated or
overlapping supports and either sign are allowed, and $N$ is finite but
otherwise unrestricted. The continuous feasible set is the full graph
$\Gamma_h=\{h(x):x\in[-1,1]^n\}$. Retaining originals makes its Boolean
witnesses $h(w)$ distinct. Let $\Phi\in\R[z]_{\le2r}$ satisfy the exact
polynomial identity
\begin{equation}\label{eq:lift-objective-pullback}
 \Phi(h(x))=F(x).
\end{equation}
Agreement merely on Boolean points is not the hypothesis here.

For $B=\prod_{j=1}^N S'_j\subseteq[-1,1]^N$, each $S'_j$ nonempty,
use the local oracle \eqref{eq:xor-node-preorder}--\eqref{eq:xor-node-quadratic}
in the $z$ coordinates, with objective $\Phi$. In addition grant
\begin{equation}\label{eq:lift-all-identities}
 L[qv]=0\quad\hbox{if }q(h(x))\equiv0,
 \qquad\deg q+\deg v\le2r.
\end{equation}
Thus every polynomial graph identity through the available degree, with
every allowed multiplier, can be imposed. Granting all of them only
strengthens the lower bound; a particular generating set may be used
instead.
A polynomial vanishing on the entire continuous graph has identically
zero pullback because the original cube has nonempty interior.
The quadratic requirement is
\begin{equation}\label{eq:lift-box-hull}
 (L[z],L[zz^T])\in\conv\{(z,zz^T):z\in B\},
\end{equation}
the hull of the entire coordinate domain, not the hull of
$\Gamma_h\cap B$, which would already optimize a linear or quadratic
objective exactly. These quadratic inequalities enter linearly on moments;
the oracle does not close them under arbitrary products or localizers.

\begin{theorem}[Bounded monomial transfer, including order one]
\label{thm:monomial-transfer}
Let $r,D\ge1$ be integers with $rD\ge2$. Suppose a Boolean
pseudoexpectation $\mathcal E$ is available through degree $4rD$, is
positive on squares through degree $2rD$, satisfies
$\mathcal E[\chi_e]=b_e$ for every clause, and has every character moment
through its available degree in $\{0,-1,1\}$. For any map
\eqref{eq:signed-monomial-map} and objective
\eqref{eq:lift-objective-pullback}, every finite cover of $\Gamma_h$ by
coordinate products certified by the preceding order-$r$ oracle at
$T_*>0$ contains at least
\begin{equation}\label{eq:monomial-transfer-count}
 2^{mT_*/(\Delta D)}
\end{equation}
regions. Arbitrary nonclosed coordinate sets and arbitrary numbers of
repeated or overlapping supports are allowed.
\end{theorem}
\begin{proof}
Let $B$ contain $h(w)$ for $w\in\pmcube$. Put
\begin{equation}\label{eq:lift-restricted-supports}
 R=\{j:\{-1,1\}\not\subseteq S'_j\},\qquad
 C=\bigcup_{j\in R}A_j,\qquad U=[n]\setminus C,
 \qquad \widetilde h(x_U)=h(w_C,x_U).
\end{equation}
Each coordinate in $R$ is now fixed to its witnessed sign. An unrestricted
coordinate becomes a signed character on $U$, possibly constant, and both
of its endpoints belong to its coordinate domain. Set
\begin{equation}\label{eq:lift-node-functional}
 L[p(z)]=\mathcal E[p(\widetilde h(x_U))].
\end{equation}
Pullback multiplies degree by at most $D$, so $L$ is defined through
lifted degree $2r$ and is normalized. As before, this is deterministic
substitution followed by marginalization, with no conditioning step.

We verify all localizers with an explicit degree budget. Suppose
$\sum_j\deg g_j+2\deg P\le2r$, where every $g_j$ is a locally valid
univariate polynomial in a lifted coordinate. A factor on a restricted
coordinate becomes a nonnegative constant. For any other coordinate,
Boolean reduction of $g_j$ after pullback equals
\[
 g_j(1)\frac{1+\widetilde h_{k_j}}{2}
 +g_j(-1)\frac{1-\widetilde h_{k_j}}{2},
 \qquad g_j(1),g_j(-1)\ge0,
\]
where $k_j$ is the coordinate used by factor $g_j$; these indices may
repeat. A constant factor needs no indicator and is removed first. Expanding the
nonconstant factors expresses the product as a sum of nonnegative
coefficients times products $I$ of parity indicators. Overlapping supports
cause no difficulty: all these factors commute and are idempotent in the
Boolean quotient, so $I^2=I$, even when $I=0$. With
$a=\sum_j\deg g_j$ and $d=\deg P$,
\begin{equation}\label{eq:lift-degree-accounting}
 \deg I\le Da,\qquad \deg\widehat P\le Dd,
 \qquad\widehat P=P\circ\widetilde h,
 \qquad\deg(I\widehat P)\le D(a+d)\le2rD.
\end{equation}
Consequently $I\widehat P^2=(I\widehat P)^2$ modulo Boolean equations,
within degree $4rD$. Original square positivity proves the localizer.
This includes every repeated bound-slack product, not just individual
bound multipliers. If a univariate equality vanishes on its
coordinate set, it vanishes at the fixed value or at both permitted
endpoints; its pulled-back Boolean reduction times any allowed multiplier
is zero. All graph identities in \eqref{eq:lift-all-identities} vanish
as literal polynomials after pullback and substitution. Their multiplied
pullbacks have degree at most $2rD$, so every stated equality is respected.

To verify \eqref{eq:lift-box-hull}, form the augmented first/second moment
matrix of $L$. A lifted linear polynomial pulls back to degree at most
$D\le2rD$, so this matrix is positive semidefinite. Every diagonal entry
is one because $\widetilde h_j^2=1$ after Boolean reduction. Every other
entry is an original character moment, multiplied by a fixed sign, and
therefore belongs to $\{0,-1,1\}$. Lemma~\ref{lem:signed-moment-law}
provides a finite Boolean distribution on the lifted coordinates with
exactly this matrix. For $j\in R$, its coordinate mean equals the fixed
sign $h_j(w)$. A Boolean variable with mean $1$ or $-1$ equals that sign
almost surely. All other coordinates have both endpoints in their sets.
Thus the realizing distribution is supported in $B$, including when the
sets are nonclosed. It may leave $\Gamma_h$: the graph moment identities
have already been proved for $L$, and no common graph-supported
distribution realizing its higher moments is asserted.

Finally, \eqref{eq:lift-objective-pullback} gives
$L[\Phi]=\mathcal E[F(w_C,x_U)]$. Only clauses meeting $C$ change their
costs. Their number is at most $\Delta|C|$. As in the unlifted proof,
each substituted clause is $(1-\sigma\chi_A)/2$ with $|A|\le3$.
The squares $(1\pm\chi_A)^2$ have degree at most six; they are available
because $2rD\ge4\ge3$. Therefore all such pseudo-costs lie in $[0,1]$,
and
\begin{equation}\label{eq:lift-cost-restriction}
 \operatorname{LB}_r(B)\le L[\Phi]\le\frac{\Delta|C|}{m}.
\end{equation}
Certification forces $|C|\ge mT_*/\Delta$.

For witness counting, write each original Boolean sign as $w_i=(-1)^{t_i}$
with $t_i\in\mathbb F_2$. Each restricted coordinate imposes one equation
with row $\ind_{A_j}$ and a right-hand side determined by $\sigma_j$ and
its permitted endpoint. Let $A_R$ be this binary row matrix and
$q=\rank_{\mathbb F_2}A_R$. The system is consistent, since it contains
$w$, and has exactly $2^{n-q}$ solutions. Every such solution meets all
endpoint restrictions and all unrestricted sets contain both endpoints,
so the witness fraction of $B$ is exactly $2^{-q}$.

Select $q$ basis rows from the original rows of $A_R$. If a column is zero
in each basis row, it is zero in their entire span, hence zero in every
row. The union of all supports is therefore the union of the basis-row
supports. Each basis support has at most $D$ elements, giving
\begin{equation}\label{eq:parity-rank-support}
 |C|\le Dq,\qquad
 2^{-q}\le2^{-|C|/D}\le2^{-mT_*/(\Delta D)}.
\end{equation}
The union bound over the graph witnesses proves the theorem. Counting
restricted lifted coordinates alone would not prove this bound: repeated
parities can create arbitrarily many restrictions with rank one.
\end{proof}

The proof needs $r\ge1$, $rD\ge2$, and $\Phi$ of degree at most $2r$,
but not $r\ge2$. The unlifted cubic objective still requires $r\ge2$
because the node functional must evaluate it. The all-character assumption is stronger
than quadratic realizability in Theorem~\ref{thm:xor-transfer}; a lifted
pair of coordinates can represent a character on as many as $2D$ original
variables. Lemma~\ref{lem:xor-width-moments} supplies precisely the required
signed moments whenever $4rD\le w$.

\begin{corollary}[A quadratic formulation with a linear objective]
\label{cor:xor-quadratic-formulation}
For every sufficiently large original dimension $n$, the family of
Theorem~\ref{thm:xor-family} has an equivalent continuous formulation
with $N=n+2m\le17n$ coordinates, $2m$ quadratic equalities, and a linear
objective. At every integer order $r\ge1$ with $12r\le an$, certificates
at absolute tolerance $1/16$ or relative tolerance $1/2$ require at least
\begin{equation}\label{eq:quadratic-formulation-count}
 2^{7n/3072}\ge2^{7N/52224}
\end{equation}
coordinate regions under the oracle of Theorem~\ref{thm:monomial-transfer}.
This includes branching on all pair and clause coordinates at order one.
\end{corollary}
\begin{proof}
Order the three distinct variables of each clause as $(i,j,k)$ and introduce
\begin{equation}\label{eq:xor-factorable-formulation}
 u_e=x_ix_j,\qquad v_e=u_ex_k,\qquad
 \Phi(x,u,v)=\frac1m\sum_e\frac{1-b_ev_e}{2}.
\end{equation}
All coordinates have bounds $[-1,1]$. For every original $x$, these
equalities determine exactly $u_e=x_ix_j$ and $v_e=x_ix_jx_k$, which lie
in the stated bounds. Conversely every feasible tuple has those values.
Thus the feasible set is the continuous monomial graph, the objectives
agree as polynomials after substitution, and the degree parameter is
$D=3$. Repeated clauses can have separate auxiliaries; neither support
repetitions nor auxiliary occurrence counts alter the theorem. The source
condition is $4rD=12r\le an$, which allows $r=1$ for all sufficiently
large $n$. Using $T_*\ge1/16$, $m\ge7n$, and $\Delta\le64$ in
\eqref{eq:monomial-transfer-count} gives the first bound. The identity
$N=n+2m\le17n$ gives the second.
\end{proof}

\begin{corollary}[Support degree versus region count]\label{cor:degree-regions}
For the same family, whenever $r\ge1$, $rD\ge2$, and $4rD\le an$,
the fixed-tolerance lower bound for degree-$D$ monomial coordinates is
\begin{equation}\label{eq:degree-region-tradeoff}
 2^{7n/(1024D)}.
\end{equation}
Within this degree regime, a number of regions polynomial in original
dimension $n$ requires $D=\Omega(n/\log n)$. For fixed $D$, the lower
bound is exponential in $n$ even with arbitrarily many auxiliaries and
node order up to a sufficiently small constant multiple of $n$.
\end{corollary}
\begin{proof}
Apply Theorem~\ref{thm:monomial-transfer} with the family parameters. If
the number of regions is at most $Cn^k$ for constants $C,k$, taking
base-two logarithms in \eqref{eq:degree-region-tradeoff} gives
$7n/(1024D)\le\log_2 C+k\log_2 n$. This is the asserted necessary
condition. It is not a sufficient construction and makes no assertion
outside $4rD\le an$.
\end{proof}
The exponent is in the original dimension $n$, not automatically in an
arbitrarily large $N$; likewise ``polynomial'' in this corollary means
polynomial in $n$. The endpoint interpolation of
Theorem~\ref{thm:local-graph-lift} handles functions of one original
coordinate at a time, and its proof does not remove the auxiliary-degree
loss for these multivariable monomials. Here literal pullback and the
parity-rank argument give a different, explicit $4rD$ degree budget.
